%%%%%%%%%%%%%%%%%%%%%%%%%%%%%%%%%%%%%%%%%%%%%%%%%%%%%%%%%%%%%%%%%
% GIÁO ÁN STEM: ECOFOOTPRINT - TÍNH DẤU CHÂN CARBON
% Môn: Toán 10/11 | Chủ đề: Hàm số bậc nhất & Thống kê
%%%%%%%%%%%%%%%%%%%%%%%%%%%%%%%%%%%%%%%%%%%%%%%%%%%%%%%%%%%%%%%%%
\documentclass[12pt,a4paper]{article}

% ============= GÓI CƠ BẢN =============
\usepackage[utf8]{inputenc}
\usepackage[T1]{fontenc}
\usepackage[vietnamese]{babel}
\usepackage{graphicx}
\usepackage{amsmath,amssymb,amsthm}
\usepackage{array}
\usepackage{booktabs}
\usepackage{multirow}
\usepackage{tabularx}
\usepackage{colortbl}
\usepackage{tikz}
\usetikzlibrary{shapes,arrows,positioning,calc,shadows}
\usepackage{pgfplots}
\pgfplotsset{compat=1.18}

% ============= FONT =============
\usepackage{newtxtext, newtxmath}
\usepackage[scaled=0.92]{helvet}
\usepackage{microtype}

% ============= MÀU SẮC (THEME XANH LÁ) =============
\usepackage[svgnames]{xcolor}
\definecolor{primary}{RGB}{34,139,34}      % ForestGreen
\definecolor{primarylight}{RGB}{144,238,144} % LightGreen
\definecolor{accent}{RGB}{255,140,0}       % DarkOrange
\definecolor{mathblue}{RGB}{0,51,102}      % NavyBlue
\definecolor{bgsoft}{RGB}{240,255,240}     % Honeydew
\definecolor{TableHeader}{gray}{0.9}

% ============= TRANG =============
\usepackage[left=2cm,right=2cm,top=2.5cm,bottom=2cm]{geometry}
\usepackage{fancyhdr}
\pagestyle{fancy}
\fancyhf{}
\renewcommand{\headrule}{\color{primary}\hrule width\textwidth height 1.5pt}
\fancyhead[L]{\color{primary}\sffamily\textbf{DỰ ÁN STEM: ECOFOOTPRINT}}
\fancyhead[R]{\color{primary}\sffamily\textbf{TOÁN HỌC VÌ MÔI TRƯỜNG}}
\fancyfoot[C]{\sffamily\thepage}

% ============= GÓI HỖ TRỢ =============
\usepackage{enumitem}
\usepackage{tcolorbox}
\tcbuselibrary{skins, breakable}
\usepackage{hyperref}
\hypersetup{colorlinks=true, linkcolor=primary, urlcolor=accent}

% ============= TCOLORBOX CUSTOM =============
\newtcolorbox{sectionbox}[1]{
    enhanced, breakable, colback=white, colframe=primary, boxrule=1pt,
    fonttitle=\bfseries\sffamily\Large, coltitle=white, colbacktitle=primary,
    title={#1},
    shadow={2mm}{-1mm}{0mm}{black!20},
    before upper={\phantomsection\addcontentsline{toc}{section}{#1}}
}

\newtcolorbox{mathbox}[1][]{
    enhanced, breakable, 
    colback=bgsoft, colframe=primary, boxrule=1pt, arc=3mm,
    title={\sffamily\bfseries\color{white} [TOÁN] KIẾN THỨC TOÁN HỌC #1},
    colbacktitle=primary, fonttitle=\bfseries\sffamily
}

\newtcolorbox{activitybox}[1]{
    enhanced, breakable,
    colback=white, colframe=accent, boxrule=1pt, arc=3mm,
    title={\sffamily\bfseries\color{white} [HOẠT ĐỘNG] #1},
    colbacktitle=accent, fonttitle=\bfseries\sffamily
}

% ============= BẮT ĐẦU TÀI LIỆU =============
\begin{document}

% ========== TRANG BÌA ==========
\begin{titlepage}
\begin{tikzpicture}[remember picture, overlay]
\fill[primary] (current page.south west) rectangle ($(current page.south east)+(0,1.5)$);
\fill[primary] (current page.north west) rectangle ($(current page.north east)+(0,-1.5)$);
\node[yshift=-1cm] at (current page.north) {\color{white}\Huge\bfseries\sffamily GIÁO ÁN STEM DỰ THI};
\end{tikzpicture}

\vspace{3cm}
\begin{center}
    {\Huge\bfseries\sffamily\color{primary} ECOFOOTPRINT}\\[1cm]
    {\LARGE\bfseries\sffamily CÔNG CỤ TÍNH DẤU CHÂN CARBON CÁ NHÂN}\\[0.5cm]
    \rule{0.6\textwidth}{2pt}\\[1cm]
    {\Large\itshape\sffamily Ứng dụng Toán học trong đo lường phát thải và bảo vệ môi trường}\\[2cm]

    \begin{tcolorbox}[width=0.7\textwidth, colback=bgsoft, colframe=primary, arc=5mm]
        \centering
        \large\sffamily
        \begin{tabular}{rl}
            \textbf{Môn học:} & Toán học (Kết hợp Sinh học, GDCD) \\
            \textbf{Khối lớp:} & 10 - 11 \\
            \textbf{Chủ đề Toán:} & Thống kê & Hàm số bậc nhất \\
            \textbf{Thời lượng:} & 2 tiết (90 phút) \\
        \end{tabular}
    \end{tcolorbox}
    
    \vfill
    {\large\sffamily TP. Hồ Chí Minh, Năm 2026}
\end{center}
\end{titlepage}

\tableofcontents
\newpage

% ========== NỘI DUNG ==========

\begin{sectionbox}{I. TỔNG QUAN CHỦ ĐỀ}
\phantomsection\subsection*{1. Mục tiêu và Ý nghĩa}\addcontentsline{toc}{subsection}{1. Mục tiêu và Ý nghĩa}
\textbf{Vấn đề thực tiễn:} Biến đổi khí hậu đang là thách thức toàn cầu. Việt Nam cam kết đạt phát thải dòng bằng 0 (Net Zero) vào năm 2050. Tuy nhiên, khái niệm "Dấu chân Carbon" (Carbon Footprint) vẫn còn trừu tượng với học sinh.

\textbf{Giải pháp STEM:} Sử dụng Toán học để \textit{định lượng} tác động môi trường của mỗi cá nhân. Học sinh sẽ xây dựng công thức tính lượng $CO_2$ phát thải từ các hoạt động thường ngày (điện, nước, đi lại) và lập trình công cụ tính toán tự động.

\phantomsection\subsection*{2. Kiến thức tích hợp}\addcontentsline{toc}{subsection}{2. Kiến thức tích hợp}
\begin{itemize}
    \item \textbf{Toán học (M):} 
    \begin{itemize}
        \item \textbf{Hàm số bậc nhất ($y=ax+b$):} Mô hình hóa mối quan hệ giữa lượng tiêu thụ và khí thải. Ví dụ: $E = k \cdot x$ (k là hệ số phát thải).
        \item \textbf{Thống kê:} Thu thập số liệu hóa đơn, tính trung bình, vẽ biểu đồ so sánh.
    \end{itemize}
    \item \textbf{Công nghệ (T):} Sử dụng ứng dụng web \texttt{EcoFootprint} để tính toán và trực quan hóa dữ liệu.
    \item \textbf{Khoa học (S):} Hiểu về hiệu ứng nhà kính, chu trình Carbon.
    \item \textbf{Xã hội (S):} Ý thức trách nhiệm công dân toàn cầu.
\end{itemize}
\end{sectionbox}

\begin{sectionbox}{II. TIẾN TRÌNH DẠY HỌC (5E MODEL)}

\phantomsection\subsection*{1. ENGAGE (Gắn kết - 10 phút)}\addcontentsline{toc}{subsection}{1. ENGAGE (Gắn kết - 10 phút)}
\begin{activitybox}{Dấu chân vô hình}
\textbf{Hoạt động:}
\begin{enumerate}
    \item GV chiếu video ngắn: "Một ngày của bạn thải ra bao nhiêu khí CO2?".
    \item GV đặt câu hỏi: \textit{"Em đi học bằng xe máy 5km, em bật điều hòa 8 tiếng. Làm sao biết hành động nào hại môi trường hơn?"}
    \item HS thảo luận và nhận ra nhu cầu cần một \textbf{công cụ đo lường} chính xác.
\end{enumerate}
\end{activitybox}

\phantomsection\subsection*{2. EXPLORE (Khám phá - 15 phút)}\addcontentsline{toc}{subsection}{2. EXPLORE (Khám phá - 15 phút)}
\begin{activitybox}{Thu thập số liệu thô}
HS làm việc theo nhóm, sử dụng \textbf{Phiếu học tập 1} để liệt kê các nguồn phát thải:
\begin{itemize}
    \item Điện năng (kWh/tháng): Xem hóa đơn điện gia đình.
    \item Nước sạch ($m^3$/tháng): Xem hóa đơn nước.
    \item Di chuyển (km/tuần): Đo quãng đường đi học.
    \item Rác thải (kg/ngày): Ước lượng.
\end{itemize}
\textbf{Câu hỏi định hướng:} \textit{"Làm sao đổi từ kWh điện sang kg $CO_2$? Cần một hệ số quy đổi nào?"}
\end{activitybox}

\phantomsection\subsection*{3. EXPLAIN (Giải thích - 20 phút)}\addcontentsline{toc}{subsection}{3. EXPLAIN (Giải thích - 20 phút)}
\begin{mathbox}[: Hàm số tính phát thải]
Để tính lượng khí thải $E$ (kg $CO_2$) từ lượng tiêu thụ $x$ (đơn vị tiêu thụ), ta dùng mô hình hàm số bậc nhất:
$$E = EF \times x$$
Trong đó \textbf{EF (Emission Factor)} là hệ số phát thải (theo số liệu Bộ TN\&MT):
\begin{itemize}
    \item \textbf{Điện:} $EF_{dien} \approx 0.72$ kg $CO_2$/kWh (nhiệt điện than chiếm tỷ trọng lớn).
    \item \textbf{Xăng (Xe máy):} $EF_{xe} \approx 0.05$ kg $CO_2$/km.
    \item \textbf{Nước sạch:} $EF_{nuoc} \approx 0.3$ kg $CO_2$/$m^3$.
    \item \textbf{Rác hữu cơ:} $EF_{rac} \approx 1.5$ kg $CO_2$/kg (do sinh khí Metan).
\end{itemize}
Tổng dấu chân Carbon:
$$E_{total} = (0.72 \cdot x_{dien}) + (0.05 \cdot x_{xe}) + \dots$$
\end{mathbox}

\phantomsection\subsection*{4. ELABORATE (Áp dụng - 30 phút)}\addcontentsline{toc}{subsection}{4. ELABORATE (Áp dụng - 30 phút)}
\begin{activitybox}{Tính toán với EcoFootprint App}
HS truy cập ứng dụng web: \href{https://stem-1h8.pages.dev/eco.html}{\texttt{stem-1h8.pages.dev/eco.html}} (Dự kiến deploy).
\begin{enumerate}
    \item Nhập số liệu cá nhân đã thu thập vào App.
    \item App tự động tính toán và vẽ biểu đồ tròn (cơ cấu phát thải).
    \item HS so sánh kết quả của mình với:
    \begin{itemize}
        \item Mức trung bình Việt Nam: $\approx 2.3$ tấn/người/năm.
        \item Mục tiêu toàn cầu: $< 2.0$ tấn/người/năm.
    \end{itemize}
    \item HS đề xuất "Kịch bản giảm thải": Nếu giảm dùng máy lạnh 1h/ngày thì tiết kiệm được bao nhiêu kg $CO_2$ một năm?
\end{enumerate}
\end{activitybox}

\phantomsection\subsection*{5. EVALUATE (Đánh giá - 15 phút)}\addcontentsline{toc}{subsection}{5. EVALUATE (Đánh giá - 15 phút)}
\begin{activitybox}{Cam kết xanh}
HS trình bày báo cáo nhóm:
- Tổng lượng phát thải ước tính.
- Nguồn phát thải lớn nhất (thường là điện hoặc xe máy).
- Cam kết hành động cụ thể (VD: Đi xe đạp 2 ngày/tuần).
GV đánh giá dựa trên Rubric (đính kèm).
\end{activitybox}

\end{sectionbox}

\newpage
\begin{sectionbox}{III. PHỤ LỤC}
\phantomsection\subsection*{Phụ lục 1: Bảng tra cứu Hệ số phát thải (EF)}\addcontentsline{toc}{subsection}{Phụ lục 1: Bảng tra cứu Hệ số phát thải (EF)}
(Nguồn tham khảo: IPCC \& Bộ TN\&MT)
\begin{center}
\begin{tabular}{|l|c|l|}
\hline
\rowcolor{primary!20} \textbf{Nguồn} & \textbf{Hệ số (kg $CO_2$)} & \textbf{Ghi chú} \\
\hline
Điện lưới & 0.7208 / kWh & Cao do nhiệt điện \\
Xe máy & 0.050 / km & Xe 110cc \\
Ô tô con & 0.180 / km & Xe 5 chỗ xăng \\
Xe buýt & 0.025 / km & Tính theo đầu người \\
Nước sạch & 0.300 / $m^3$ & Xử lý và bơm \\
Thịt bò & 27.0 / kg & Chăn nuôi phát thải cao \\
Rau củ & 0.50 / kg & Canh tác \\
\hline
\end{tabular}
\end{center}

\phantomsection\subsection*{Phụ lục 2: Rubric Đánh giá}\addcontentsline{toc}{subsection}{Phụ lục 2: Rubric Đánh giá}
\begin{tabularx}{\textwidth}{|>{\bfseries}l|X|X|X|}
\hline
\rowcolor{primary!20} Tiêu chí & Mức 1 (Cần cố gắng) & Mức 2 (Đạt) & Mức 3 (Tốt) \\
\hline
Thu thập dữ liệu & Số liệu ước lượng, thiếu chính xác & Có số liệu từ hóa đơn, nhưng chưa đầy đủ & Số liệu chi tiết, nguồn gốc rõ ràng (hóa đơn, đo đạc) \\
\hline
Tính toán Toán học & Sai công thức hoặc kết quả & Áp dụng đúng công thức $E=EF \cdot x$ & Tính đúng, hiểu ý nghĩa con số, biết quy đổi đơn vị \\
\hline
Sử dụng Công nghệ & Nhập liệu sai vào App & Nhập đúng, ra kết quả & Phân tích được biểu đồ, so sánh kịch bản trên App \\
\hline
Đề xuất giải pháp & Chung chung (bảo vệ MT) & Có giải pháp cụ thể & Giải pháp định lượng được hiệu quả (giảm bao nhiêu kg) \\
\hline
\end{tabularx}

\end{sectionbox}

\end{document}
